\documentclass[../ICIT-Paper-2024.tex]{subfiles}
\usepackage{array}
\usepackage{color}
\usepackage[table]{xcolor}
\usepackage{multirow}
\usepackage{multicol}
\usepackage{multirow}
\usepackage{hyperref} 
\usepackage{tablefootnote}
\begin{document}
\paragraph{}

Applying the credit scoring model of credit scores, the data will be divided into 5 credit levels (table~\ref{tab3}). However, labeling the data with crypto wallets encounters a major issue. The problem is that we cannot determine which level each wallet truly falls into, as there is no clear credit evaluation mechanism for each level. Even experts cannot establish clear rules to define the 5 credit levels. For example, a wallet may have a balance of \$1 trillion and be liquidated once in the previous month. This wallet could potentially fall into level 2 due to the liquidation event, or into level 4 due to its large balance. These characteristics resemble those of Fuzzy Data.

% Áp dụng mô hình chấm điểm tín dụng của credit score, dữ liệu sẽ được chia thành 5 mức tín dụng (Bảng 1). Tuy nhiên, việc gán nhãn dữ liệu với tập ví điện tử gặp phải một vấn đề lớn. Vấn đề đó là chúng ta không thể biết ví thực sự rơi vào thang nào. Do không có cơ chế đánh giá tín dụng rõ ràng cho từng mức. Ngay cả chuyên gia cũng không thể có cơ sở nào rõ ràng để xây dựng tập luật xác định 5 mức tín dụng. Ví dụ, ta có ví có balance 1 tr đô và bị thanh lý 1 lần trong 1 tháng trước. Ví này sẽ có thể rơi vào thang 2 vì đặc điểm bị thanh lý và rơi vào thang 4 vì lượng tiền ví này lớn. Những đặc điểm này tương đồng với đặc điểm của Fuzzy Data.
\begin{table}[h!]
\caption{FICO Credit score range}
\label{tab:credit_score_ranges}
\centering
\begin{tabular}{|c|c|m{7cm}|}
\hline
\textbf{Credit Score Ranges} & \textbf{Rating} & \textbf{Description} \\ 
\hline
300-579 & Poor & This credit score is well below the average score of U.S. consumers and demonstrates to lenders that the borrower may be a risk. \\ 
\hline
580-669 & Fair & This credit score is below the average score of U.S. consumers, though many lenders will approve loans with this score. \\ 
\hline
670-739 & Good & Likely to qualify for credit at competitive rates \\ 
\hline
740-799 & Very Good & Eligible for favorable rates and terms \\ 
\hline
800-850 & Excellent & Typically qualifies for the best rates and terms available \\ 
\hline 
\end{tabular}
\end{table}

According to ~\citeauthor{viertl2011statistical}~\cite{viertl2011statistical}, Fuzzy data are all kinds of data that cannot be represented as precise numbers or cannot be precisely classified. The analysis of Fuzzy Data has been researched and compiled in many books, such as ~\citeauthor{bandemer2012fuzzy}~\cite{bandemer2012fuzzy}, ~\citeauthor{viertl2011statistical}~\cite{viertl2011statistical}, and ~\citeauthor{nguyen2006fundamentals}~\cite{nguyen2006fundamentals}, etc. Machine learning algorithms that solve problems related to Fuzzy data or logic can be easily found on Google Scholar.

% Theo sách, Fuzzy data are all kinds of data which cannot be represented as precise numbers or cannot be precisely classified. Việc phân tích các Fuzzy Data đã được nghiên cứu và tổng hợp tại nhiều cuốn sách như book1, book2, và, book3, etc. Các giải thuật học máy giải quyết các bài toán liên quan đến Fuzzy data hay Fuzzy logic có thể tìm thấy dễ dàng trên gg scholar. 


~\citeauthor{denoeux2011maximum}~\cite{denoeux2011maximum} introduces a method for estimating parameters in parametric statistical models with fuzzy observations linked to underlying crisp random samples. The method maximizes the observed-data likelihood, which is the probability of the fuzzy data, using the EM algorithm, enabling solutions to various statistical problems involving fuzzy data. ~\citeauthor{jane2019review}~\cite{jane2019review}reviews how two techniques (machine learning and fuzzy logic) facilitate knowledge-based decision-making. This paper details their efficiency and the steps involved. ~\citeauthor{yang1996class}~\cite{yang1996class} introduces a class of fuzzy clustering procedures specifically designed for fuzzy data, as opposed to traditional fuzzy clustering techniques that handle crisp data with fuzzy set theory-based class memberships. The newly derived procedures, called fuzzy c-numbers (FCN) clusterings, are tailored to manage different types of fuzzy data. The paper details the construction of these FCNs for LR-type, triangular, trapezoidal, and normal fuzzy numbers.
% Một số bài báo điển hình như (1) bài báo 1 introduces a method for estimating parameters in parametric statistical models with fuzzy observations linked to underlying crisp random samples. The method maximizes the observed-data likelihood, which is the probability of the fuzzy data, using the EM algorithm, enabling solutions to various statistical problems involving fuzzy data. (2) bài báo 2 reviews how two techniques (machine learning and fuzzy logic) facilitate knowledge-based decision making. This paper details their efficiency and the steps involved. (3) bài báo 3 introduces a class of fuzzy clustering procedures specifically designed for fuzzy data, as opposed to traditional fuzzy clustering techniques that handle crisp data with fuzzy set theory-based class memberships. The newly derived procedures, called fuzzy c-numbers (FCN) clusterings, are tailored to manage different types of fuzzy data. The paper details the construction of these FCNs for LR-type, triangular, trapezoidal, and normal fuzzy numbers.

We decided to label the set of electronic wallets using Fuzzy Data. In other words, our labels are Fuzzy Labels. Our labeling idea is inspired by Fuzzy c-means clustering~\cite{bezdek1984fcm}. Fuzzy clustering is a powerful unsupervised method for data analysis and model construction. In many situations, fuzzy clustering is more natural than hard clustering. Objects on the boundaries between several classes are not forced to fully belong to one of the classes but are instead assigned membership degrees between 0 and 1, indicating their partial membership.

% Chúng tôi quyết định gán nhãn có tập ví điện tử theo dạng Fuzzy Data. Nói cách khác labels của chúng tôi là các Fuzzy Labels. Ý tưởng gán nhãn của chúng tôi được lấy cảm hứng từ Fuzzy c-mearns clustering. Fuzzy clustering is a powerful unsupervised method for the analysis of data and construction of models. In many situations, fuzzy clustering is more natural than hard clustering. Objects on the boundaries between several classes are not forced to fully belong to one of the classes, but rather are assigned membership degrees between 0 and 1 indicating their partial membership. 
The labeling is carried out in detail as table~\ref{tab4}:


\begin{table}[h!]
\caption{Label description by score range}\label{tab4}
\centering
\begin{tabular}{|p{2.5cm}|p{2.5cm}|p{2.5cm}|p{2.5cm}|p{2.5cm}|}
\hline
\textbf{Poor} & \textbf{Fair} & \textbf{Good}& \textbf{Very Good}& \textbf{Excellent} \\

\hline
Accounts with assets under \$\numprint{1000} have been liquidated more than three times. & Accounts with assets less than \$\numprint{1000} have been liquidated fewer than three times and have deposits. & Accounts with assets over \$\numprint{1000} and a borrow-to-asset ratio of 30\%-40\% has not been liquidated. & Accounts with assets over \$1 million and a borrow-to-asset ratio of 20\%-30\% has not been liquidated. & Accounts with assets greater than \$1 million with a borrow-to-asset ratio below 20\% has no history of asset liquidation. \\
\hline
\multicolumn{2}{|p{5cm}|}{
    Accounts with a borrow-to-asset ratio over 40\% have been liquidated more than once. \newline
    Accounts with assets under \$\numprint{100000} and a debt ratio over 60\% have been liquidated. \newline
    Accounts with assets over \$\numprint{1000} and a borrow-to-asset ratio below 40\% have been liquidated at least once. \newline
    Accounts with assets under \$\numprint{1000} have been liquidated fewer than three times and have no deposits.
} & \multicolumn{2}{p{5cm}|}{
    Accounts with assets ranging from \$\numprint{1000} to \$\numprint{10000} have not been liquidated.
} & \cellcolor{gray!30} \\
\hline
\cellcolor{gray!30} & \multicolumn{2}{|p{5cm}|}{
    Accounts with assets under \$\numprint{100000} and a debt ratio from 40\% to 60\% have not been liquidated.
} & \multicolumn{2}{p{5cm}|}{
    Accounts with assets ranging from \$\numprint{10000}to \$\numprint{100000} and a borrow-to-asset ratio below 20\% have not been liquidated.\newline
    Accounts with assets over \$\numprint{100000} have no borrowings, have deposits and have never been liquidated.
} \\
\hline
\end{tabular}
\end{table}

Currently, these parameters are being selected based on our analysis of 6 million positively interacting wallets across 7 EVM chains over the past year. In the future, we will further investigate improving the mechanism for selecting these thresholds. Based on this mechanism, an automated analysis and labeling system will be deployed. This labeling system will serve as the foundation for developing and evaluating prediction models, allowing us to better direct and refine classification algorithms based on clear and quantifiable financial characteristics of the researched subjects.

\end{document}